% generated by stats/analyze_sweeps.py -- do not edit by hand
% \input this, or copy the blocks you want into the chapter.

\begin{figure}[t]
\centering
\includegraphics[width=\textwidth]{assets/fig1_mechanism_effects.pdf}
\caption[Effect of each coordination mechanism on episode success, by deployment geometry]{Effect of each coordination mechanism on episode success, by deployment geometry. Paired exact McNemar within each cell; significance after Benjamini--Hochberg correction over all cells. Blue = the mechanism helps, red = it hurts. Rows of panels separate the two target modes: the \texttt{tail} row is the rare-target regime coordination is meant to pay for. The $N=2$ and $N=4$ bars are separate paired experiments, so their difference is unpaired and is shown rather than tested.\label{fig:fig1_mechanism_effects}}
\end{figure}

\begin{figure}[t]
\centering
\includegraphics[width=\textwidth]{assets/fig2_success_matrix.pdf}
\caption[Episode success across the deployment cells]{Episode success across the deployment cells. Deployment communication radius as trained. The grey rule is the random-policy floor for that cell; the shaded column is the training condition, and every other cell is out of distribution.\label{fig:fig2_success_matrix}}
\end{figure}

\begin{figure}[t]
\centering
\includegraphics[width=\textwidth]{assets/fig3_comms_range.pdf}
\caption[Restricting the communication range at deployment]{Restricting the communication range at deployment. Faint lines are individual deployment cells, bold is the mean over them. Every policy was trained with an unlimited link, so any change along this axis is a train/deploy mismatch rather than a learned response.\label{fig:fig3_comms_range}}
\end{figure}

\begin{figure}[t]
\centering
\includegraphics[width=\textwidth]{assets/fig4_success_at_budget.pdf}
\caption[Success within a step budget]{Success within a step budget. Pooled over the deployment cells. Episodes never solved count as failures at every budget, so the curves are free of survivorship bias. The normalized area under a curve is the time-weighted success of Table~\ref{tab:time_penalized}.\label{fig:fig4_success_at_budget}}
\end{figure}

\begin{figure}[t]
\centering
\includegraphics[width=\textwidth]{assets/fig5_time_to_first_success.pdf}
\caption[Time to first success: the conditional view disagrees with the penalized one]{Time to first success: the conditional view disagrees with the penalized one. Left: median over each arm's own successes, which is survivorship-biased -- an arm that solves fewer episodes drops its slow ones and can look faster ($n$ = episodes solved, printed on each bar). Right: every episode counted, with unsolved episodes charged the full budget.\label{fig:fig5_time_to_first_success}}
\end{figure}

\begin{figure}[t]
\centering
\includegraphics[width=\textwidth]{assets/fig8_spawn_effect.pdf}
\caption[Where the swarm starts: effect of the spawn geometry on whether the target is found, and on how long the first agent takes to find it]{Where the swarm starts: effect of the spawn geometry on whether the target is found, and on how long the first agent takes to find it. Spawn regimes are ordered by how dispersed the swarm is at $t=0$: \texttt{origin} places every agent on a single point at the corner of the coast, \texttt{random} is the uniform draw the policies were trained on (shaded), \texttt{max\_dist} spaces them evenly along the coastline. Faint lines are the two target modes separately, bold is the pooled result. Both rows describe the FIRST arrival only: an episode ends when any one agent reaches the zone. Top: episode success rate. Bottom: penalized time $E[\min(t,T)]$, which charges an unsolved episode the full budget and therefore cannot be improved by failing more often -- the conditional median, which can, is in Table~\ref{tab:spawn_effect}. The grey random-policy line is the control: where it moves in step with the others, the spawn regime changed the difficulty of the task itself rather than the value of what the policies learned.\label{fig:fig8_spawn_effect}}
\end{figure}

\begin{figure}[t]
\centering
\includegraphics[width=\textwidth]{assets/fig9_hard_targets.pdf}
\caption[Cost of a rare target: how much later, and how much less often, the swarm arrives]{Cost of a rare target: how much later, and how much less often, the swarm arrives. Empirical distribution of the time to first arrival, pooled over the three spawn geometries. The value each curve reaches at the right-hand edge is that arm's success rate and the shape of its rise is its speed, so both quantities are read off one plot; episodes that are never solved simply never fire, which is why no curve reaches~1. Note that the curves are still climbing when the budget cuts them off at \texttt{max\_steps}: none of these arms has finished finding what it would eventually find, so every success rate quoted here is a rate \emph{within the budget} rather than an asymptote. The tick on each curve marks the median over that arm's own successes -- the conditional number of Table~\ref{tab:hard_targets}, shown against the distribution it conditions away. Left: an ordinary target anywhere in the domain. Right: a target drawn from the rarest few per cent of the salinity distribution in its depth plane, which shrinks the success zone. Axes are shared, so the cost of rarity is the rightward shift plus the drop in plateau.\label{fig:fig9_hard_targets}}
\end{figure}

\begin{figure}[t]
\centering
\includegraphics[width=\textwidth]{assets/fig10_scaling_n.pdf}
\caption[Effect of swarm size on finding the target and on how long it takes]{Effect of swarm size on finding the target and on how long it takes. Pooled over all six deployment cells, 600 episodes per point. Left: episode success rate, counting an episode once whichever agent arrives first. Right: penalized time to first arrival $E[\min(t,T)]$, which charges an unsolved episode the full budget and so cannot be improved by failing more often. The $N=2$ and $N=4$ episodes are separate draws from the same task distribution rather than matched pairs -- \texttt{reset()} places the agents before it draws the target, so a seed gives a different task at a different swarm size -- and these are therefore unpaired comparisons of rates. Table~\ref{tab:scaling_n} adds the search-overlap and travel cost of the extra agents, the margin over an untrained swarm of the same size, and the comparison against what non-interacting agents would already achieve.\label{fig:fig10_scaling_n}}
\end{figure}

\begin{figure}[t]
\centering
\includegraphics[width=\textwidth]{assets/fig6_arrival_by_rank_N2.pdf}
\caption[Arrival of the $r$-th agent, not just the first]{Arrival of the $r$-th agent, not just the first ($N = 2$). Empirical arrival CDF of the $r$-th agent over all episodes; the plateau is the fraction that ever arrived within the budget. Censoring is administrative at \texttt{max\_steps}, so this equals the Kaplan--Meier estimate. Plateau heights are bounded by the previous rank -- an $r$-th agent cannot arrive unless $r-1$ already did -- so the informative comparison is between arms within a panel.\label{fig:fig6_arrival_by_rank_N2}}
\end{figure}

\begin{figure}[t]
\centering
\includegraphics[width=\textwidth]{assets/fig6_arrival_by_rank_N4.pdf}
\caption[Arrival of the $r$-th agent, not just the first]{Arrival of the $r$-th agent, not just the first ($N = 4$). Empirical arrival CDF of the $r$-th agent over all episodes; the plateau is the fraction that ever arrived within the budget. Censoring is administrative at \texttt{max\_steps}, so this equals the Kaplan--Meier estimate. Plateau heights are bounded by the previous rank -- an $r$-th agent cannot arrive unless $r-1$ already did -- so the informative comparison is between arms within a panel.\label{fig:fig6_arrival_by_rank_N4}}
\end{figure}
